% =====================================================================
% trmreuse.tex -- machinery for re-using Castalia TRM LaTeX fragments
% inside this Beamer deck, unmodified and un-copied.
% Originally written for docs/publications/castalia_lab_talk; this copy
% differs only in the path block below (relative, repo-anchored).
%
% The TRM's figure fragments (pgfplots plots, circuitikz schematics,
% generated TikZ block diagrams and booktabs tables) are \input directly
% from the TRM source tree, so a regenerated TRM updates this deck.  Two
% things have to be arranged for that to work:
%
%   1. path + macro context -- the fragments assume the TRM's cwd and a
%      handful of the TRM's own \newcommands.  Both are supplied below.
%   2. float suppression + captions -- a TRM caption is 3-5 sentences of
%      body text and belongs on a page, not a slide.  \trmfig strips the
%      float wrapper, the caption and the label, and scales what is left.
% =====================================================================

% --- where the shared sources live --------------------------------------
% This deck sits in implementations/asic/castalia/docs/deck/, next to the
% published TRM.pdf, and every path below is RELATIVE to that directory, so
% the deck moves with the repository.
%
%   \AnalogSrc  implementations/asic/castalia/analog/   the hand-written analog
%               chapter and the maestro2tex fig_/tab_/sch_ fragments -- the
%               tracked source that LatexUserGuide.CopyAnalogChapter copies
%               into the TRM build.  Read from the source, not the copy.
%   \TRMinc     platform/common/latex/TRM/include/      the GENERATED system
%               diagrams and defines.tex.  Gitignored; run `make generate`
%               (or `make chip`) in platform/common first.
%   \Assets     assets/web/                             the layout renders.
%   \ISCASroot  the ISCAS'27 paper (outside the repo) -- the concept figure
%               only, guarded with \IfFileExists so the deck builds without it.
\newcommand{\VestaRoot}{../../../../../}
\newcommand{\AnalogSrc}{../../analog/}
\newcommand{\TRMroot}{\VestaRoot platform/common/latex/TRM/}
\newcommand{\TRMinc}{\TRMroot include/}
\newcommand{\Assets}{\VestaRoot assets/web/}
\providecommand{\MaestroRoot}{\AnalogSrc}
\newcommand{\ISCASroot}{/home/mseminario2/work/ieee/ISCAS27/latek/}
\makeatletter\edef\input@path{{\TRMroot}{\TRMinc}{\AnalogSrc}{\ISCASroot}}\makeatother
\graphicspath{{\Assets}{\TRMroot figures/}{\ISCASroot Figures/}}
\pgfplotsset{table/search path={\ISCASroot}}

% --- the TRM macros the fragments use --------------------------------
% Verbatim from platform/common/latex/packages-commands.tex (all of them
% are plain \texttt wrappers) plus the generated defines.
\providecommand{\pin}[1]{\texttt{#1}}
\providecommand{\peripheral}[1]{\texttt{#1}}
\providecommand{\register}[1]{\texttt{#1}}
\providecommand{\bitfield}[1]{\texttt{#1}}
\providecommand{\asminline}[1]{\texttt{#1}}
\providecommand{\cinline}[1]{\texttt{#1}}
\providecommand{\forthinline}[1]{\texttt{#1}}
\providecommand{\bashinline}[1]{\texttt{#1}}
\input{\TRMinc defines.tex}

% \memsection -- structural, not a stub: AddressSpaceDiagram needs the real
% one.  Verbatim from platform/common/latex/packages-commands.tex line 93.
\providecommand{\memsection}[4]{%
	\bytefieldsetup{bitheight=#3\baselineskip}%
	\bitbox[]{10}{%
		\texttt{#1}
		\\
		\vspace{#3\baselineskip}
		\vspace{-2\baselineskip}
		\vspace{-#3pt}
		\texttt{#2}
	}%
	\bitbox{16}{#4}%
}

% --- the TRM figure theme (v* styles) --------------------------------
% Verbatim from platform/common/latex/packages-commands.template.tex: the
% generated diagrams (PowerDomainDiagram, ...) draw with these.
\definecolor{vestaRed}{rgb}{0.75,0,0}
\definecolor{vestaRedText}{rgb}{0.70,0,0}
\definecolor{vestaInk}{rgb}{0,0,0}
\tikzset{
	% Boxes. Three weights of fill, one stroke, one corner radius. The block
	% diagrams draw with vbox: a white box with the thin grey outline, the
	% RP2040 datasheet look; the filled weights stay defined for the figures
	% that tell two kinds of node apart by weight.
	vbox/.style={line width=0.6pt, rounded corners=2pt, draw=black!70, fill=white},
	vblock/.style={thick, rounded corners=2pt, draw=vestaInk, fill=black!8},
	vblocklt/.style={thick, rounded corners=2pt, draw=vestaInk, fill=black!3},
	vblockw/.style={thick, rounded corners=2pt, draw=vestaInk, fill=white},
	vblockem/.style={line width=1.1pt, rounded corners=2pt, draw=vestaInk, fill=black!12},
	% The bus bar, and the package boundary.
	vbar/.style={thick, draw=vestaInk, fill=black!15},
	vbound/.style={draw=vestaRed, line width=1.2pt},
	% A grouping region: a thin dashed outline that leaves the paper white, never a
	% full-bleed grey band.
	vregion/.style={draw=black!35, line width=0.5pt, dashed, rounded corners=3pt},
	% Text: sans throughout, in two sizes only. vname (\small) names a block, vlab
	% (\footnotesize) carries everything secondary, and nothing is set below
	% \footnotesize, so the smallest label is 8 pt at natural size. The older names
	% stay defined, at the same floor, for the emitters that still use them.
	vname/.style={font=\sffamily\small, align=center, inner sep=1pt},
	vlab/.style={font=\sffamily\footnotesize, align=center, inner sep=1pt},
	vhd/.style={font=\sffamily\footnotesize, align=center, inner sep=1pt, anchor=north},
	vbc/.style={font=\sffamily\footnotesize, align=center, inner sep=1pt},
	vsm/.style={font=\sffamily\footnotesize, align=center, inner sep=1pt},
	vnote/.style={font=\sffamily\footnotesize\itshape, text=black!55, align=left, inner sep=1.5pt},
	vgroup/.style={font=\sffamily\small\itshape, text=black!55, align=left, fill=white, inner sep=1.5pt},
	vredlab/.style={font=\sffamily\small\bfseries, text=vestaRedText, align=left},
	% Lines: one arrow head for the whole manual, Stealth, on every style that carries
	% a head, and a junction dot where a wire really joins another.
	vwire/.style={semithick, draw=vestaInk},
	vflow/.style={->, >={Stealth[length=2mm]}, semithick, draw=vestaInk},
	vbus/.style={<->, >={Stealth[length=2mm]}, thick, draw=vestaInk},
	vaccent/.style={->, >={Stealth[length=2mm]}, semithick, draw=vestaRed},
	vghost/.style={dashed, draw=black!45, semithick},
	vdot/.style={circle, fill=vestaInk, inner sep=0pt, minimum size=3pt},
	% A mux: a trapezium standing on its narrow side, inputs on the wide (left) edge,
	% the select code beside each input.
	vmux/.style={trapezium, trapezium angle=70, shape border rotate=270, line width=0.6pt, draw=black!70, fill=white, inner sep=1pt},
}

% --- schematics: heavier strokes for projection ------------------------
% The TRM's circuitikz schematics are drawn for print at 0.4-0.6 pt.  On a
% slide they vanish, so every path inside a \trmfit'd fragment is drawn one
% step heavier; explicit "thick" rails stay relatively heavier still.
\newcommand{\trmslidestrokes}{%
  \tikzset{every path/.append style={line width=0.9pt}}%
  \ctikzset{bipoles/thickness=1.3}%
}

% \memrow / \memgap -- the address-column rows AddressSpaceDiagram uses since
% the 2026-09 TRM (verbatim from packages-commands.template.tex).
\providecommand{\memrow}[3]{%
	\bytefieldsetup{bitheight=#2\baselineskip}%
	\bitbox[]{10}{%
		\texttt{#1}
		\\
		\vspace{#2\baselineskip}
		\vspace{-2\baselineskip}
		\vspace{-#2pt}
		\mbox{}
	}%
	\bitbox{16}{\sffamily\small #3}%
}
\providecommand{\memgap}[2]{%
	\bytefieldsetup{bitheight=#2\baselineskip}%
	\bitbox[]{10}{%
		\texttt{#1}
		\\
		\vspace{#2\baselineskip}
		\vspace{-2\baselineskip}
		\vspace{-#2pt}
		\mbox{}
	}%
	\bitbox{16}[bgcolor=black!8]{}%
}

% --- the maestro2tex categorical palette -----------------------------
% Slots 1-5 of the validated reference palette, in documented order; the
% block masters declare these with \providecolor, individual fig_*.tex do not.
\definecolor{mtxA}{HTML}{2A78D6}
\definecolor{mtxB}{HTML}{EB6834}
\definecolor{mtxC}{HTML}{1BAF7A}
\definecolor{mtxD}{HTML}{EDA100}
\definecolor{mtxE}{HTML}{E87BA4}

% --- \trmfig[scale]{fragment.tex} ------------------------------------
% Inputs a TRM fragment with its float, caption and label removed and the
% result scaled.  Everything is done inside a group so the redefinitions
% never escape onto the rest of the slide.
\newsavebox{\trmbox}
\newcommand{\trmfig}[2][0.62]{%
  \begingroup
    \renewenvironment{figure}[1][]{\ignorespaces}{}%
    \renewenvironment{table}[1][]{\ignorespaces}{}%
    \renewcommand{\caption}[2][]{}%
    \renewcommand{\label}[1]{}%
    \renewcommand{\FloatBarrier}{}%
    \sbox{\trmbox}{\begin{minipage}{\linewidth}\centering\input{#2}\end{minipage}}%
    \scalebox{#1}{\usebox{\trmbox}}%
  \endgroup}

% \trmtab -- same, for a booktabs table fragment (no scaling by default).
\newcommand{\trmtab}[2][1.0]{\trmfig[#1]{#2}}

% \trmfigw[width fraction]{fragment.tex} -- same, but fitted to a fraction
% of \linewidth instead of scaled by a factor.  Use this for the wide
% generated block diagrams, whose natural width is 20-34 cm.
\newcommand{\trmfigw}[2][1.0]{%
  \begingroup
    \renewenvironment{figure}[1][]{\ignorespaces}{}%
    \renewenvironment{table}[1][]{\ignorespaces}{}%
    \renewcommand{\caption}[2][]{}%
    \renewcommand{\label}[1]{}%
    \renewcommand{\FloatBarrier}{}%
    \sbox{\trmbox}{\input{#2}}%
    \resizebox{#1\linewidth}{!}{\usebox{\trmbox}}%
  \endgroup}

% \trmfigh{height}{fragment.tex} -- fitted to a stated height.  The usable
% body height of a 16:9 Berkeley frame is about 6.6 cm.
\newcommand{\trmfigh}[2]{%
  \begingroup
    \renewenvironment{figure}[1][]{\ignorespaces}{}%
    \renewenvironment{table}[1][]{\ignorespaces}{}%
    \renewcommand{\caption}[2][]{}%
    \renewcommand{\label}[1]{}%
    \renewcommand{\FloatBarrier}{}%
    \sbox{\trmbox}{\input{#2}}%
    \resizebox*{!}{#1}{\usebox{\trmbox}}%
  \endgroup}

% =====================================================================
% \duckph{width}{what the missing artwork is}
% Placeholder for artwork that does not exist yet -- the top-level GDS
% plot, die photographs, bonding diagrams.  A tikzducks duck in a dashed
% frame: impossible to mistake for a real figure on a projector, and
% impossible to leave in by accident.
% =====================================================================
\newcommand{\duckph}[3][graduate, book, bookcolour=red!60!black]{%
  \begin{center}
  \setlength{\fboxsep}{6pt}%
  \tikz[baseline]{\node[draw=red!70!black, very thick, dashed, rounded corners=3pt,
        inner sep=8pt, align=center, text width=#2]{%
     {\bfseries\color{red!70!black} ARTWORK PENDING}\\[4pt]
     \begin{tikzpicture}[scale=0.95]
       \duck[#1]
     \end{tikzpicture}\\[4pt]
     {\small #3}};}
  \end{center}}

% =====================================================================
% \trmfit[natural width]{max width}{max height}{fragment.tex}
%
% Fits a fragment inside a BOX, scaling uniformly by whichever limit binds
% first, so it can be dropped into a column without overflowing either way.
% This is the one to reach for by default.
%
% WHY THE OPTIONAL FIRST ARGUMENT MATTERS.  Every generated pgfplots
% fragment declares "width=0.82\linewidth, height=6.2cm" -- a RELATIVE
% width against an ABSOLUTE height.  Input inside a half-slide beamer
% column, \linewidth is about 5.9 cm, so the axis comes out 4.8 cm wide by
% 6.2 cm tall: portrait, and every trace squashed in x.  Scaling that box
% afterwards cannot undo it -- the aspect is already wrong.
%
% So \trmfit sets \hsize/\linewidth/\columnwidth to a wide value BEFORE the
% \input, which is what the fragment's 0.82\linewidth then resolves
% against.  The plot is laid out landscape and only then scaled down to
% fit.  15 cm gives roughly 12.3 x 6.2 cm, a 2:1 axis.  Pass a smaller
% value for a fragment that should come out squarer, larger for a wider one.
% Circuitikz fragments (which use \resizebox{\linewidth}{!}) are unaffected:
% they are laid out at the natural width and scaled back uniformly.
% =====================================================================
\newlength{\trmnat}
\newcommand{\trmfit}[4][15cm]{%
  \begingroup
    \renewenvironment{figure}[1][]{\ignorespaces}{}%
    \renewenvironment{table}[1][]{\ignorespaces}{}%
    \renewcommand{\caption}[2][]{}%
    \renewcommand{\label}[1]{}%
    \renewcommand{\FloatBarrier}{}%
    \setlength{\trmnat}{#1}%
    % \linewidth is set as a VALUE for the fragment to resolve against, but
    % the box is NOT wrapped in a minipage of that width -- a minipage would
    % pin \wd\trmbox to \trmnat and the fit below would then scale a 33 cm
    % TikZ diagram as though it were 15 cm wide, and it would run off the slide.
    \sbox{\trmbox}{\hsize=\trmnat \linewidth=\trmnat \columnwidth=\trmnat
                   \trmslidestrokes\input{#4}}%
    % natural TOTAL height (ht+dp: a tabular is centred on the baseline, so dp
    % is half of it), then the height the box would have if scaled to width #2
    \dimen0=\ht\trmbox \advance\dimen0 by \dp\trmbox
    \divide\dimen0 by 64                      % keep the product inside 2^31 sp
    \count0=\wd\trmbox  \divide\count0 by 65536
    \count2=\dimexpr#2\relax \divide\count2 by 65536
    \ifnum\count0<1 \count0=1 \fi
    \multiply\dimen0 by \count2
    \divide\dimen0 by \count0
    \multiply\dimen0 by 64
    \ifdim\dimen0>\dimexpr#3\relax
      \resizebox*{!}{#3}{\usebox{\trmbox}}%    height binds (starred: TOTAL height,
                                               % a tabular sits centred on the baseline)
    \else
      \resizebox{#2}{!}{\usebox{\trmbox}}%     width binds
    \fi
  \endgroup}
